Put \(N=M+1\). For either arm \(c\), the elementary equivalence \(\lceil x\rceil\le n\) if and only if \(x\le n\), for integral \(n\), gives
\[
k_c\le M\iff N(1-\varepsilon_c)\le N-1\iff\varepsilon_c\ge\frac1N.
\tag{1}
\]

Under \(\texttt{def:disjoint-orbit-law}\), the arm \(a\) potential outcome equals \(K\) on target-supported arm \(a\) orbits and zero elsewhere, while every arm \(1-a\) potential outcome is zero. These coordinates are constant on every \(\mathcal G\)-orbit. Complete the prototype with any fixed covariate array and let \(G\) be uniform on the finite group, as the definition prescribes. For fixed \(g\in\mathcal G\), the product \(gG\) is uniform, so \(g(GL^*)\) and \(GL^*\) have the same law. The constructed law therefore satisfies finite-group latent invariance. Its potential outcomes lie in \(\{0,K\}\subseteq[-K,K]\), so \(\texttt{def:invariant-laws}\) places it in \(\mathcal P_{\mathcal G,K}\). Independent copies, the known assignment law, and the stipulated auxiliary selector draws furnish the required replicated experiment.

If \(q^{\mathrm{obs}}_{w,a}(O)>0\), the disjoint-support premise says that \(O\) is not target-supported, and the construction assigns outcome zero there. Since the observable selector chooses only realized arm \(a\) members and the predictor is identically zero, every selected arm \(a\) calibration score is zero. All arm \(1-a\) outcomes and its predictor are also zero, so every calibration score in that arm is zero. Consequently, for \(c\in\{0,1\}\),
\[
r_{c,(k_c)}=0\quad\text{if }k_c\le M,\qquad r_{c,(k_c)}=+\infty\quad\text{if }k_c=M+1.
\tag{2}
\]
If both ranks are finite, clipping leaves \(I^{\boldsymbol\varepsilon}_{0,w}=I^{\boldsymbol\varepsilon}_{1,w}=\{0\}\), and hence \(\mathcal C^{\boldsymbol\varepsilon}_w=\{0\}\). Target independence puts the selected target in a target-supported orbit almost surely, so
\[
\tau_T=\begin{cases}K,&a=1,\\-K,&a=0.\end{cases}
\tag{3}
\]
Because \(K>0\), (3) is not in \(\{0\}\), proving zero conditional coverage.

If exactly one rank is at the boundary, (2) and \(\texttt{prop:asymmetric-rank-boundary}\) give one interval \([-K,K]\) and the other interval \(\{0\}\). In either order their Minkowski difference is \([-K,K]\), which contains both endpoint values in (3), so this witness has coverage one. If both ranks are at the boundary, the same proposition gives \(\mathcal C^{\boldsymbol\varepsilon}_w=[-2K,2K]\), again containing (3). Equation (1) proves the final rank equivalence. On the symmetric diagonal the two ranks coincide; the only regimes are therefore the both-finite zero-coverage case and the both-boundary coverage-one case, which is exactly the symmetric witness calculation.